% =====================================================================
% KidneyMTL - manuscript draft
% Numbers marked \todo{} are produced by scripts/aggregate_results.py etc.
% (see README "Analyse"). Tables are \input directly from results/tables/.
% Before submission: switch to the target journal's template, verify every
% reference, fill every \todo{}, and remove this header.
% =====================================================================
\documentclass[11pt]{article}
\usepackage[margin=1in]{geometry}
\usepackage{amsmath,amssymb,booktabs,graphicx,hyperref,xcolor,multirow}
\usepackage[numbers,sort&compress]{natbib}
\newcommand{\todo}[1]{\textcolor{red}{[TODO: #1]}}
\graphicspath{{../results/figures/}{../results/explain/}{figures/}}

\title{Beyond Leaky Benchmarks: Leakage-Aware Multi-Task Learning for\\
Kidney Disease Classification and Stone Localisation in CT}
\author{Utkarsh Kumar\thanks{\todo{affiliation, e-mail}}}
\date{}

\begin{document}
\maketitle

% ---------------------------------------------------------------------
\begin{abstract}
Deep-learning studies on public kidney CT datasets routinely report accuracies above 99\%. We show that a
large part of this performance is produced by the evaluation protocol rather than by the models. On the
widely used CT-KIDNEY dataset (12,446 slices; normal, cyst, tumour, stone), 99.3\% of test images in a random
image-level split have a near-duplicate (perceptual-hash distance $\le$4/64 bits) in the training set, the
dataset contains 517 sets of byte-identical files, and the common practice of augmenting before splitting
places transformed copies of 98.6\% of test images in the training data. We propose (i)~a leakage-aware
evaluation protocol that groups near-duplicate slices and performs stratified group 5-fold cross-validation,
and (ii)~KidneyMTL, a multi-task network whose shared EfficientNet-B3 backbone is trained jointly for
four-class kidney disease classification and kidney-stone localisation on two partially labelled datasets,
with a CBAM classification head, an anchor-free CenterNet detection head and uncertainty-based task weighting.
Under the leakage-free protocol KidneyMTL reaches a macro-F1 of \todo{XX.X $\pm$ X.X}\% and a stone
detection mAP@0.5 of \todo{XX.X}\%, outperforming \todo{N} ImageNet baselines and a re-implementation of the
recent CDR-VGG19 model (\todo{XX.X}\%; McNemar $p<$\todo{}). Re-training the same models under the random and
augment-then-split protocols inflates macro-F1 by up to \todo{XX} points, which explains the near-perfect
results reported previously. Code, folds and all trained configurations are publicly available.
\end{abstract}

\noindent\textbf{Keywords:} kidney stone; CT; multi-task learning; data leakage; object detection;
explainable AI; calibration

% ---------------------------------------------------------------------
\section{Introduction}
Kidney stone disease affects a large and increasing share of the population, and non-contrast CT is the
reference imaging modality for its diagnosis. Cysts and renal tumours appear on the same scans, so an
automated reader should both separate these conditions and, for stones, show \emph{where} the stone is.
Most recent deep-learning studies frame the problem as image-level classification on public Kaggle datasets
and report accuracies between 98\% and 99.9\%~\citep{islam2022vision,devi2025cdr,yildirim2021deep}.

Two methodological issues make these numbers hard to interpret. First, the public CT-KIDNEY dataset
\citep{islam2022vision} has no patient identifiers and consists of consecutive, highly similar slices from the
same patients, so a random image-level split places near-copies of each test image in the training set.
Second, several studies augment the dataset (e.g., 21 variants per image in~\citep{devi2025cdr}) \emph{before}
splitting it, so rotated or flipped versions of the same image appear in both training and test data. Both are
instances of the leakage problems that have been identified as a major cause of irreproducible results in
machine-learning-based science and medical imaging~\citep{kapoor2023leakage,roberts2021pitfalls,varoquaux2022failures}.

A second limitation is that almost all prior models output only a class label. A detector that localises the
stone is more useful clinically and gives a direct check on whether the classifier relies on the lesion.
Kidney stone datasets with box annotations are, however, small and contain no non-stone classes, while large
classification datasets contain no boxes.

This paper addresses both points. Our contributions are:
\begin{enumerate}
  \item A \textbf{leakage audit} of CT-KIDNEY based on perceptual hashing, and a \textbf{grouped, stratified
  5-fold protocol} that keeps every group of near-duplicate slices inside a single fold (Section~\ref{sec:protocol}).
  \item A \textbf{controlled leakage experiment} that re-trains identical models under three protocols
  (grouped, random, augment-then-split) and quantifies how much each protocol inflates reported performance
  (Section~\ref{sec:leakage}).
  \item \textbf{KidneyMTL}, a multi-task network trained on two \emph{partially labelled} datasets: four-class
  classification from CT-KIDNEY and stone localisation from a box-annotated stone dataset. It uses CBAM
  attention, a CenterNet head on an FPN neck, and homoscedastic-uncertainty task weighting (Section~\ref{sec:method}).
  \item A \textbf{tri-channel input} $[\,\text{CT},\ \text{CLAHE(CT)},\ \text{contour map}\,]$ that keeps the
  contour prior of~\citep{devi2025cdr} without discarding intensity information.
  \item A \textbf{complete evaluation}: seven baselines including a re-implementation of
  CDR-VGG19~\citep{devi2025cdr}, twelve ablations, paired significance tests, calibration analysis,
  quantitative explainability, an external cross-dataset test and an efficiency analysis.
\end{enumerate}

% ---------------------------------------------------------------------
\section{Related work}
\paragraph{Kidney stone and kidney disease classification.}
CNN classifiers for kidney stones on coronal CT include sequential CNNs~\citep{yildirim2021deep} and ensembles;
\citet{islam2022vision} released CT-KIDNEY and compared vision transformers with CNNs (Swin reaching
99.3\% accuracy under a random split). \citet{devi2025cdr} combine a contour-detected input with a VGG-19
backbone and a normalised residual block (CDR-VGG19), reporting 99.61\% accuracy after 21$\times$
augmentation followed by an 80:10:10 split. \todo{add 3--5 recent papers on CT-KIDNEY and report their split
protocol; most use random splits}.

\paragraph{Stone detection and segmentation.}
\citet{elton2022deep} detect and segment stones on non-contrast CT with a 3D U-Net and \citet{parakh2019urinary}
evaluate stone detection and generalisation across scanners. \todo{add YOLO-based stone detectors}.

\paragraph{Multi-task learning.}
Sharing a representation across related tasks can improve generalisation~\citep{caruana1997multitask}.
\citet{kendall2018multi} weight task losses by learned homoscedastic uncertainty, which we adopt to balance
classification and detection without manual tuning.

\paragraph{Leakage and evaluation in medical imaging.}
\citet{roberts2021pitfalls} found that most COVID-19 imaging models were not fit for clinical use, partly
because of duplicated and overlapping data. \citet{kapoor2023leakage} catalogue leakage types across
disciplines, and \citet{varoquaux2022failures} discuss evaluation failures specific to medical imaging.
Reporting guidelines such as CLAIM~\citep{mongan2020claim} require patient-level separation of data subsets.

% ---------------------------------------------------------------------
\section{Datasets and evaluation protocol}\label{sec:protocol}
\paragraph{CT-KIDNEY (classification).} 12,446 axial CT slices (JPEG) labelled Normal (5,077), Cyst (3,709),
Tumor (2,283) and Stone (1,377)~\citep{islam2022vision}. No patient identifiers are provided.

\paragraph{Kidney-stone boxes (detection).} 1,300 coronal CT images with \todo{N} stone bounding boxes
(Roboflow export of the Kaggle ``kidney stone images'' dataset, CC~BY~4.0), split by its authors into
1,054/123/123 train/validation/test images. Roboflow created two augmented versions of each source image. We
verified that no source image appears in more than one split.

\paragraph{Leakage audit.} We compute a 64-bit DCT perceptual hash~\citep{zauner2010phash} of every CT-KIDNEY
slice and link slices whose Hamming distance is $\le\tau$. Connected components form \emph{groups}. Table~\ref{tab:audit}
summarises the audit. With $\tau{=}4$ the data form 229 groups (largest: 705 slices). Larger thresholds merge
almost half of the dataset into a single component through transitive chaining ($\tau{=}6$: 5,664 slices), so we
use $\tau{=}4$, the strictest threshold that still yields usable folds. In addition, 517 sets of byte-identical
files (1,034 images) exist.

\begin{table}[t]\centering
\caption{Near-duplicate audit of CT-KIDNEY (mean over 5 folds). ``NN'' = nearest training image of each test image.}
\label{tab:audit}
\begin{tabular}{lcc}\toprule
Protocol & Test images with NN $\le$4 bits & Median NN distance (bits) \\\midrule
Random image split & 99.3\% & 0.0 \\
Augment-then-split~\citep{devi2025cdr} & 98.6\%$^\dagger$ & -- \\
Grouped (ours) & 0.0\% & 10.4 \\\bottomrule
\multicolumn{3}{l}{\footnotesize $^\dagger$ test source images whose augmented copies are in the training set.}
\end{tabular}
\end{table}

\paragraph{Protocol.} Groups are distributed over five parts with StratifiedGroupKFold. In fold $k$, part $k$ is
the test set, part $(k{+}1)\bmod 5$ the validation set (model selection and early stopping) and the remaining three
parts the training set (60/20/20). Every image is therefore tested exactly once, which yields out-of-fold
predictions for the whole dataset. Augmentation is applied online and to training images only. For the leakage
experiment we also build (a) a stratified \emph{random} image-level 5-fold split and (b) the \emph{augment-then-split}
protocol of~\citep{devi2025cdr}: every image is expanded into five deterministic variants (rotations, flips,
intensity changes, noise, zoom) and the variants are split at random.

\paragraph{Limitation of the proxy.} Hash groups approximate but do not equal patients. Two different patients
may fall in one group (14 of 229 groups contain more than one class), and distant slices of the same patient may
fall in different groups. The grouped protocol is therefore a strict lower bound on leakage, not a guarantee of
patient-level independence.

% ---------------------------------------------------------------------
\section{Method}\label{sec:method}
\todo{Figure 1: architecture diagram (backbone, CBAM head, FPN, CenterNet head, two data streams).}

\paragraph{Input representation.} Each grey-level slice $g$ is mapped to three channels
$[\,g,\ \mathrm{CLAHE}(g),\ K(g)\,]$, where $K(g)$ is a multi-level iso-intensity contour map: after Gaussian
smoothing, contours are traced at the Otsu threshold and at the 85th, 95th and 99th intensity percentiles of the
foreground and drawn with level-dependent intensity. Hyper-dense structures such as stones therefore receive
nested high-level contours. Unlike~\citep{devi2025cdr}, the contour map is added to the image, not substituted
for it, so an ImageNet-pretrained backbone can use all three views.

\paragraph{Backbone and classification head.} An ImageNet-pretrained EfficientNet-B3~\citep{tan2019efficientnet}
provides feature maps at strides 4, 8, 16 and 32. The stride-32 map passes through CBAM~\citep{woo2018cbam},
global average pooling, dropout, a 256-unit layer with SiLU, and a linear layer with four outputs.

\paragraph{Detection head.} A top-down FPN fuses the four feature levels into a 128-channel stride-4 map,
followed by a CenterNet head~\citep{zhou2019objects} that predicts a stone-centre heatmap, the box size and a
sub-pixel offset. Training targets are Gaussian peaks with the CenterNet radius. Boxes are decoded by
$3{\times}3$ max-pooling non-maximum suppression and the top-50 peaks.

\paragraph{Losses.} The classification loss is the focal loss~\citep{lin2017focal} ($\gamma{=}2$) with
inverse-square-root class weights and label smoothing 0.05. The detection loss is
$\mathcal{L}_{det}=\mathcal{L}_{hm}+\mathcal{L}_{off}+2\,\mathcal{L}_{CIoU}$ with the penalty-reduced focal
heatmap loss, an $L_1$ offset loss and the CIoU loss~\citep{zheng2020diou} on boxes decoded at ground-truth centres.
The tasks are combined with learned homoscedastic uncertainty~\citep{kendall2018multi}:
$\mathcal{L}=e^{-s_c}\mathcal{L}_{cls}+s_c+e^{-s_d}\mathcal{L}_{det}+s_d$.

\paragraph{Partially labelled training.} Each optimisation step draws one CT-KIDNEY batch (classification loss
only) and one stone batch (detection loss only) and updates the shared backbone with the combined loss.
No image needs both kinds of label.

\paragraph{Implementation.} AdamW (backbone LR $2{\times}10^{-4}$, heads $10^{-3}$, weight decay $10^{-4}$),
2-epoch linear warm-up and cosine decay, at most 40 epochs, early stopping on validation macro-F1 (patience 8),
mixed precision, gradient clipping at 5, and an exponential moving average of the weights (decay 0.999) used for
evaluation. Inputs are $320^2$ (classification) and $512^2$ letterboxed (detection). Training augmentation:
horizontal flip, affine ($\pm15^\circ$, $\pm15\%$ scale, $\pm6\%$ shift), brightness/contrast, gamma, Gaussian
noise, blur and coarse dropout. Test-time augmentation: horizontal-flip averaging. Hardware: \todo{GPU};
software versions: \todo{torch, timm}. All hyperparameters are listed in \texttt{configs/default.yaml}.

\paragraph{Baselines.} ResNet-50, DenseNet-121, EfficientNet-B3, ConvNeXt-T, ViT-S/16 and VGG-19-BN, fine-tuned
from ImageNet with cross-entropy and the same augmentation, schedule and protocol. CDR-VGG19~\citep{devi2025cdr}
is re-implemented from the paper's description (VGG-19, two 4096-unit FC layers, a 4096-320-320-4096 normalised
residual block, contour-detected $224^2$ input; 146M parameters, matching the reported size).

\paragraph{Metrics and statistics.} Accuracy, balanced accuracy, macro precision/recall/F1, macro one-vs-rest AUC,
Cohen's $\kappa$, MCC, stone-vs-rest sensitivity and specificity, expected calibration error
(15 bins)~\citep{guo2017calibration} and Brier score. We report the mean $\pm$ standard deviation over test folds
and, on the pooled out-of-fold predictions, 95\% percentile bootstrap confidence intervals. The bootstrap
resamples whole near-duplicate groups (2,000 replicates), so correlated slices are not counted as independent
evidence. Models are compared with exact McNemar tests on paired predictions, with Holm correction within each
table. Detection is evaluated with COCO-style AP at IoU 0.5 and 0.5:0.95 (101-point interpolation), together
with precision, recall and F1 at a score threshold of 0.3.

% ---------------------------------------------------------------------
\section{Results}
\subsection{Comparison with baselines}
Table~\ref{tab:main} \todo{describe: best model, margin over strongest baseline and over CDR-VGG19, per-class
behaviour (stone vs tumour confusion), significance from Table~\ref{tab:significance}}.
\IfFileExists{../results/tables/main_comparison.tex}{\input{../results/tables/main_comparison.tex}}{\todo{run aggregate\_results.py}}

\subsection{How much does the protocol inflate performance?}\label{sec:leakage}
Table~\ref{tab:leakage} and Figure~\ref{fig:leakage} \todo{report the gap between grouped and random / augment-then-split
for each model; note that CDR-VGG19 under the reference protocol approaches the published 99.61\% while it drops to
XX\% under the grouped protocol}.
\IfFileExists{../results/tables/leakage.tex}{\input{../results/tables/leakage.tex}}{}
\begin{figure}[t]\centering
\IfFileExists{../results/figures/leakage.png}{\includegraphics[width=0.7\linewidth]{leakage.png}}{\todo{figure}}
\caption{Test macro-F1 of the same architectures under three evaluation protocols.}\label{fig:leakage}
\end{figure}

\subsection{Ablation study}
\IfFileExists{../results/tables/ablation.tex}{\input{../results/tables/ablation.tex}}{}
\todo{discuss: effect of the detection task on classification (a\_no\_det), input representation (gray / CLAHE /
contour-only vs tri), attention, loss, task weighting, backbone}.

\subsection{Stone localisation}
\IfFileExists{../results/tables/detection.tex}{\input{../results/tables/detection.tex}}{}
\todo{compare KidneyMTL detection with the detection-only model (does the classification data help detection?) and
YOLOv8 as an external reference; show results/runs/kidneymtl\_main/fold0/plots/detections.png}.

\subsection{Statistical significance, calibration and efficiency}
\IfFileExists{../results/tables/significance.tex}{\input{../results/tables/significance.tex}}{}
\IfFileExists{../results/tables/efficiency.tex}{\input{../results/tables/efficiency.tex}}{}
\todo{ECE before/after; reliability diagram; params/GFLOPs vs CDR-VGG19 (146M)}.

\subsection{Explainability}
\todo{Grad-CAM / Grad-CAM++ figure (results/explain/gradcam\_comparison.png). Pointing game on the Dataset-1 test
stones: KidneyMTL XX\% vs baselines XX\% (results/explain/pointing\_game.csv), and XX\% for the detection heatmap.}

\subsection{External validation}
\todo{Stone sensitivity of each classifier on the unseen, coronal Dataset-1 test split
(results/tables/external\_validation.csv). Interpret the cross-view domain gap honestly.}

% ---------------------------------------------------------------------
\section{Discussion}
\todo{(1) Leakage: near-perfect accuracies on CT-KIDNEY mostly reflect duplicate memorisation; recommend grouped
evaluation and release of folds. (2) What multi-task learning adds: localisation output, effect on classification
and on Grad-CAM focus. (3) Clinical relevance: stone sensitivity, calibration for triage. (4) Limitations: no
patient IDs (hash groups are a proxy), 2D slices instead of volumes, JPEG-compressed images without Hounsfield
units, single-centre public data, axial vs coronal mismatch between the two datasets, no reader study.
(5) Future work: volumetric models, multi-centre DICOM data, stone composition/size.}

\section{Conclusion}
\todo{3--4 sentences.}

\paragraph{Data and code availability.} Both datasets are public on Kaggle. The code, configuration files and the
exact fold assignments are available at \url{https://github.com/utkarshhzz/Kidney_research}.

\bibliographystyle{unsrtnat}
\bibliography{references}
\end{document}
